\section{Records: dados como dados}
\label{sec:records}

O \citetitle{doc:jep395} tornou os records finais no JDK~16, depois de
duas rodadas de \termo{preview}{Recurso entregue para experimentação, que
exige a opção \texttt{-{}-enable-preview} e pode mudar antes de ficar
final.} (JDK~14 e 15). Um record declara um portador transparente de dados imutáveis:
o construtor, os \termo{accessors}{Métodos de leitura; num record, têm o
mesmo nome do componente, \texttt{x()} e não \texttt{getX()}.},
\texttt{equals}, \texttt{hashCode} e \texttt{toString} são derivados dos
componentes declarados.

\lstinputlisting[language=Java]{generated/source/RecordExemplo.java}

O que o \texttt{javac} gerou para \texttt{Ponto}:

\inputbytecode{RecordExemplo_Ponto}

\paragraph{Lendo a listagem.}
\begin{itemize}
\item \texttt{final class \ldots\ extends java.lang.Record}: o record não
  pode ser estendido e herda de \texttt{Record}, não de \texttt{Object}
  diretamente.
\item \texttt{private final int x} e \texttt{y}: os campos são
  \texttt{final}; o construtor só faz \texttt{putfield} uma vez em cada.
\item \texttt{x()} e \texttt{y()} são um \texttt{getfield} e um
  \texttt{ireturn}: não há custo escondido.
\item \texttt{toString}, \texttt{hashCode} e \texttt{equals} não têm
  corpo gerado pelo compilador: cada um é um único
  \texttt{invokedynamic}. Com \texttt{javap -v} (desafio no fim do
  capítulo) aparece o método de \termo{bootstrap}{Método chamado uma única
  vez, na primeira execução de um \texttt{invokedynamic}, para decidir que
  código aquele ponto de chamada vai executar. Ver a
  seção~\ref{sec:invokedynamic}.}
  \texttt{java.lang.runtime.ObjectMethods.bootstrap}, que recebe os nomes
  \texttt{"x;y"} e referências aos dois campos. A implementação fica no
  JDK, não no seu \texttt{.class}.
\end{itemize}

\paragraph{A decisão que isso muda.}
Para DTOs, eventos, chaves de mapa e respostas de API, o record substitui
a classe com getters, \texttt{equals} e \texttt{hashCode} escritos à mão
(ou gerados por Lombok). O \texttt{equals} deixa de ficar desatualizado
quando alguém acrescenta um campo e esquece de atualizá-lo, porque ele é
derivado da lista de componentes. Isso segue o conselho do
\citetitle{livro:bloch2017} no item 17, \enquote{Minimize mutability}.

\paragraph{O bug que o bytecode não impede.}
O \texttt{final} vale para o campo, não para o objeto apontado. Um
\texttt{record Pedido(List<Item> itens)} guarda a referência à lista
recebida: quem tiver essa lista ainda pode alterá-la, e o
\texttt{hashCode} do pedido muda junto. Se a imutabilidade importa, copie
no construtor compacto (\texttt{itens = List.copyOf(itens);}).

\subsection*{Record patterns}
O \citetitle{doc:jep440} (final no JDK~21) permite desmontar um record
no próprio padrão: \texttt{if (o instanceof Ponto(int x, int y))} testa o
tipo e extrai os componentes de uma vez. Combinado com
\texttt{switch} e sealed types (próxima seção), isso substitui boa parte
dos \texttt{Visitor} e das cadeias de \texttt{instanceof} com cast.
